In this section, we present several countermeasures aimed at 
improving the resilience of the signature scheme against the 
fault attacks described above.

We group the countermeasures into three categories. The 
first category consists of \emph{post-publication code-path 
checks}, which preserve both the signature size and the 
original tree-construction procedure. The second category 
comprises \emph{modifications to the path-construction 
algorithm}, which maintain the original signature size while 
altering the path-generation process. Finally, the third 
category considers \emph{tree-less GGM signatures}, which 
eliminate the tree structure entirely at the cost of a 
larger signature size.


\subsection{Post-publication path check}

This countermeasure runs checks on the output of \SeedTreeToPath.

A simple size check on the path is enough to be secure against fault models 1 and 3, 
since they result in path lengths of $1$ and $\text{golden path size}+1$ respectively.
The expected path size can be computed in linear time from $\calH$. 
A more thorough check would be to rebuild the tree using $\PathToSeedTree$ 
on the obtained path and to compare the obtained leaves to $\calH$ and the reference tree.

\paragraph{Verify-before-release, and why it is weaker than it looks.}
The most natural version of this countermeasure is the classical one: run the
verification algorithm on the signature before releasing it, and suppress the
output on mismatch. Against the attacks of Section~\ref{sec:app} this does
work in principle, because none of the three faulted paths deserialises to the
honest leaf set and the recomputed digest therefore differs from $d$.

It is weaker than it looks for a structural reason specific to this
implementation. In the MEDS reference code, $\SeedTreeToPath$ and
$\PathToSeedTree$ are not two functions but two modes of a single function,
\texttt{stree\_to\_path\_to\_stree}, selected by a \texttt{mode} parameter and
sharing the traversal loop of Listing~\ref{lst:path}. The instruction the
attacker glitches during signing is therefore \emph{the same instruction}, at
the same code offset, that the verification pass will execute. An attacker who
can place one glitch can place two, and a ChipWhisperer can emit both within a
single trace. The redundancy is nominal: the check re-executes the faulted
code rather than checking it independently.

This is the reason we recommend the decoupled construction of
Section~\ref{sec:decouple} over verification alone, and it generalises beyond
MEDS --- any implementation that folds $\Publish$ and $\Recon$ into shared code
inherits the same weakness. We have not measured the success rate of the
double-glitch experiment on our platform, and doing so is the most useful
single extension of this work: it would turn a structural observation into a
cost.

\begin{remark}
Sign-then-verify is also the wrong shape of defence for a second reason. As
noted in Section~\ref{sec:app}, the adversary does not need a valid signature;
it needs the device to emit a corrupted path. Any countermeasure of this
family must therefore suppress the output entirely on detection, not merely
flag it.
\end{remark}



\subsection{Path construction algorithm modification}
\label{sec:decouple}

Although the unified implementation of $\SeedTreeToPath$ and $\PathToSeedTree$ is efficient, 
its structural coupling of the flag tree construction and path extraction phases creates a vulnerability.
Decoupling these operations into independent processing stages, and introducing validation checks after each phase, 
adds a redundancy that ensures that any fault injected during the initial flag tree 
generation will cause a detectable mismatch during the subsequent path construction phase.

We implement this countermeasure using the following multi-stage verification procedure:

\begin{enumerate}
\item Construct a binary flag tree $\calF$ derived from $\calH$.
\item Verify that for every leaf position $i$ where $\calF[i] = 1$, no ancestral node in its path to the root is flagged as $0$.
\item Reconstruct the node path using the validated flag tree $\calF$ and the reference seed tree.
\item Assert that no leaf node explicitly flagged as $1$ is included in the final extracted path.
\end{enumerate}

\noindent\textbf{Complexity. }  All four verification steps scale linearly with the size of the tree, maintaining a time complexity of $\mathcal{O}(t)$.


\noindent\textbf{Impact on performances.}
We implement all three countermeasures and compare the execution times of the signature. 

The results are summarised in Table~\ref{tab:counter_1_cmp}. We write
$C_{size}$ for the path-size check, $C_{tree}$ for the full reconstruction
check, and $C_{flagtree}$ for the decoupled flag-tree construction of
Section~\ref{sec:decouple}.

\noindent\textit{Experimental Setup.} We evaluated MEDS under RIOT OS on two
platforms. The embedded target is a bare-metal Nordic nRF52840 DK (32-bit ARM
Cortex-M4 at 64 MHz, 256 KB RAM, 1 MB Flash) running the Toy parameter set; the
host baseline runs all parameter sets (MEDS9923--MEDS167717) natively on an
Intel Core i7 desktop under Ubuntu 24.04 LTS. Code was compiled with
\texttt{arm-none-eabi-gcc} (Cortex-M4) and \texttt{gcc} 13.3.0 (host).


\begin{table}[htbp]
\centering
\caption{MEDS signature time (microseconds) by parameter set. All values are the median value of 100 runs.}
\label{tab:counter_1_cmp}
\smallskip
\begin{tabular*}{\textwidth}{l@{\extracolsep{\fill}}cccc}
\toprule
\textbf{Parameter set} & \textbf{base} & $C_{size}$ & $C_{tree}$ & $C_{flagtree}$ \\
\midrule
Toy                    & 6845    & 6793      & 6794      & 6835          \\
MEDS9923               & 1609626	& 1582806	& 1590051.5	& 1585796 \\
MEDS13220              & 299950	& 275909	& 283484	& 283477.5         \\
MEDS41711              & 4847694.5	& 5200203	& 4962401.5	& 5360679.5 \\
MEDS55604              & 1280193	& 1381107	& 1321436	& 1307889          \\
MEDS134180             & 5009655.5	& 5074072.5	& 5140707.5	& 5735290 \\
MEDS167717             & 2950702	& 2951835.5	& 3076849.5	& 3051602.5  \\
\midrule
Toy (Cortex M4)        & 753566	& 737162	& 738843	& 737739
\\
\bottomrule
\end{tabular*}
\end{table}

All three checks are $\mathcal{O}(t)$ against a signing cost dominated by
matrix arithmetic over $\Fq$, so the expected overhead is small, and this is
what Table~\ref{tab:counter_1_cmp} shows: the largest observed figures are
$+7.9\%$ for $C_{size}$, $+4.3\%$ for $C_{tree}$ and $+14.5\%$ for
$C_{flagtree}$, the last on \textit{MEDS134180}.

These numbers should be read with care, and we prefer to say so rather than
overstate them. Several rows report the instrumented build as \emph{faster}
than the baseline --- $-8.0\%$ for $C_{size}$ on \textit{MEDS13220},
$-1.7\%$ on \textit{MEDS9923} --- which cannot be a real effect, since the
countermeasure only adds work. The noise floor of a whole-signature wall-clock
measurement on this platform is therefore of the same order as the overheads
we are trying to report. The correct experiment is to instrument
$\SeedTreeToPath$ alone in cycle counts and to report a distribution rather
than a median; we regard the numbers below as an upper bound on the order of
magnitude ($\le 15\%$ of signing time in the worst case, and negligible for
the shallow-tree parameter sets) rather than as precise measurements.
\todo[inline]{Re-run: per-function cycle counts, $\ge 1000$ runs, report
median with IQR. Also: Table~\ref{tab:counter_1_cmp} lists MEDS55604, which is
absent from Tables~\ref{tab:scheme-params} and~\ref{tab:medsoverhead}
(MEDS69497 there). Reconcile the parameter set actually benchmarked.}

\subsection{GGM-tree-less signature}

Alternative architectural designs could consider abandoning GGM tree structures entirely. 
A potential solution, adapting the fault attack countermeasure for LESS proposed by \cite{ZKFault} 
would require the response vector $v$ to explicitly contain the precomputed matrices 
$\tilde{A}_i$ and $\tilde{B}_i$ that are sampled from the published leaves $\sigma_i$.

\begin{equation}
v[i] = 
\begin{cases}
    (\tilde{A}_i . A_{h_i}^{-1}, B_{h_i}^{-1}.\tilde{B}_i), & \text{if } h_i > 0 \\
    (\tilde{A}_i,\, \tilde{B}_i), & \text{if } h_i = 0
\end{cases}
\end{equation}

While this design ensures that the secret matrices $(\tilde{A}_i, \tilde{B}_i)$ remain computationally 
infeasible to derive for any index $i$, it introduces a disadvantage by significantly inflating 
the overall signature size to the size \texttt{sig\_cm} shown in table~\ref{tab:medsoverhead}.

\begin{table}[htbp]
\centering
\caption{Comparison of Signature Sizes and Overhead across MEDS Parameter Sets. \texttt{sig\_cm} is given by $\ell_{\text{digest}} + t(\ell_{\mathbb{F}_q^{m \times m}} + \ell_{\mathbb{F}_q^{n \times n}}) + \ell_{\text{salt}}$
}
\label{tab:medsoverhead}
\begin{tabularx}{\textwidth}{l c >{\raggedleft\arraybackslash}X >{\raggedleft\arraybackslash}X r}
\toprule
\textbf{Parameter Set} & NIST Level & \texttt{sig} (Bytes) & \texttt{sig\_cm} (Bytes) & \textbf{Size Overhead} \\
\midrule
MEDS-9923   & I   & 9,896   & 677,440 & 6745.59\% \\
MEDS-13220  & I   & 12,976  & 112,960 & 770.53\%  \\
MEDS-41711  & III & 41,080  & 882,880 & 2049.17\% \\
MEDS-69497  & III & 54,736  & 232,384 & 324.55\%  \\
MEDS-134180 & V   & 132,424 & 475,072 & 258.75\%  \\
MEDS-167717 & V   & 165,332 & 277,152 & 67.63\%   \\
\bottomrule
\end{tabularx}
\end{table}